\chapter{绪论}
\label{chp:introduction}

\section{研究背景与意义}

图结构可以刻画实体之间的连接关系及其相互作用，因此常用于交通网络、社交网络、生物分子网络、知识图谱和软件依赖分析等场景。对图结构进行推理和计算，是图论、数据挖掘与人工智能研究中的基础问题。最短路径、连通性判定、匹配、最大团、最小顶点覆盖以及旅行商问题等任务，既有清晰的图论定义和算法基础，也对应路径规划、网络分析、资源调度和组合优化等实际应用需求\cite{cormen2009clrs,diestel2017graph,korte2018combinatorial}。

近年来，大语言模型在语义理解、步骤化推理和程序生成方面表现出较强能力。Transformer 架构、少样本学习、指令微调以及开源基础模型的发展，使语言模型逐渐从文本建模工具扩展为能够参与复杂任务求解的通用模型\cite{vaswani2017attention,brown2020language,ouyang2022instructgpt,touvron2023llama,bubeck2023sparks}。在链式思维提示、自一致性采样、树状思维等方法的推动下，语言模型在数学推理、代码生成和工具调用任务中已经展现出较好的问题分解与执行辅助能力\cite{wei2022chain,wang2023selfconsistency,yao2023tree,besta2024graph,gao2023pal}。基于这些能力，研究者开始将大语言模型用于图推理任务，尝试让模型根据自然语言问题理解图结构、选择合适算法，并生成相应的求解过程或程序。

\begin{figure}[htbp]
    \centering
    \includegraphics[width=0.9\textwidth]{figures/paper/intro.pdf}
    \caption{图推理任务整体示意图}
    \label{fig:intro_task_overview}
\end{figure}

图\ref{fig:intro_task_overview}给出了本文所研究图推理任务的基本流程：系统首先接收自然语言形式的任务描述，再围绕给定图结构执行最短路径、最大团、旅行商等不同类型的图推理任务，最终生成目标答案。与传统纯文本问答不同，这类任务的关键在于如何把自然语言目标准确映射到可执行的图算法流程上。

图推理与一般文本推理存在本质差异。图结构强调节点、边及其拓扑关系，输入信息通常由离散的结构对象组成，而非自然语言中的线性语义序列。为了适配语言模型，大多数现有方法需要将图序列化为文本，再让模型直接进行端到端推理或生成代码。随着图规模增大、边属性变复杂以及约束条件增多，这种串行化表示会显著增加模型理解负担，导致拓扑关系遗漏、约束违反和答案幻觉等问题。

针对这一挑战，本文关注的并非一般性地描述图推理任务难度，而是大语言模型在图推理求解中面临的流程耦合问题：任务理解、算法选择、知识调用、代码实现和正确性验证常被压缩在一次生成过程中，任何一个中间环节出现偏差，都可能造成最终答案错误或程序不可执行。对于图推理这类同时依赖结构理解、算法步骤和约束满足的任务，这种耦合会放大错误传播风险。

因此，本文关注大语言模型在图推理任务中的可靠求解问题，重点讨论如何在不同类型图任务之间保持较好的泛化能力。本文认为，单纯延长推理链，或直接更换能力更强的基础模型，并不能稳定避免图结构遗漏、算法选择偏差和程序执行错误。对于图推理任务而言，更需要在自然语言任务描述和可执行程序之间加入一个可检查的中间层，用来记录任务目标、算法步骤和约束条件。在此基础上，规划结果用于约束后续的知识检索与代码生成，执行反馈则用于反向定位并修正程序实现中的偏差和遗漏。基于这一思路，本文构建了面向图推理任务的多智能体协同求解框架，将规划的前向约束作用与反馈的反向修正作用结合起来。


\section{国内外研究现状}

为了深入研究和设计面向图推理任务的规划反馈多智能体协同求解方法，本节对现有国内外相关领域研究工作进行了系统归纳，主要从以下三个方面展开：大语言模型推理增强方法、图推理任务上的大模型研究进展，以及现有方法在图任务求解中的主要局限。通过对这些工作的梳理，本文说明选择规划反馈协同求解路线的原因。

\subsection{大语言模型推理方法研究}

大语言模型推理能力的提升，早期主要依赖提示形式设计和解码策略改进。Chain-of-Thought 通过让模型写出中间推理步骤，改善了其在算术推理和逻辑推理任务中的表现\cite{wei2022chain}。Scratchpad、过程验证器以及 Program of Thoughts 等研究进一步说明，将中间计算过程、程序执行过程或结果检查过程显式呈现出来，可以减少复杂推理中不容易被直接观察到的错误\cite{nye2021scratchpads,cobbe2021verifiers,chen2023pot}。Self-Consistency 不再只依赖单条推理路径，而是对多条候选推理结果进行聚合，从而降低偶然推理错误对最终答案的影响\cite{wang2023selfconsistency}。Least-to-Most Prompting 将原问题拆分为一组较小的子问题，并按照由易到难的顺序逐步求解\cite{zhou2023least}。Tree-of-Thoughts 和 Graph-of-Thoughts 则进一步把推理过程组织为可分支、可搜索的结构，而不是固定的一条线性链路\cite{yao2023tree,besta2024graph}。此外，Zero-Shot CoT、Selection-Inference、ReAct、Self-Refine、Reflexion、CRITIC、语言智能体树搜索以及过程验证类方法，分别从零样本推理、选择与推断分离、推理和行动结合、自反馈修正、工具交互纠错、搜索式行动规划以及过程监督等角度扩展了推理增强研究的范围\cite{kojima2022zeroshot,creswell2023selection,yao2023react,madaan2023selfrefine,shinn2023reflexion,critic2024,zhou2024lats,lightman2023processreward}。这些工作表明，在复杂任务中显式组织中间步骤，往往比直接生成最终答案更稳定。

除文本提示层面的改进外，另一条研究路线是将语言模型接入外部工具、执行环境或程序系统。PAL 将自然语言问题转化为程序生成与执行问题，使模型负责生成求解逻辑，而把具体计算交给解释器完成\cite{gao2023pal}。Toolformer 和 ToolLLM 更强调模型对外部 API 的识别、选择和调用能力，使语言模型可以在生成过程中主动利用外部功能\cite{schick2023toolformer,qin2023toolllm}。Gorilla 和 AlphaCode 则分别从 API 对齐和复杂程序生成角度说明，语言模型在外部环境支持下可以完成更复杂的任务\cite{gorilla2024,liang2022alphacode}。这类研究的共同特点是：模型不再被要求独立完成所有计算细节，而是通过生成程序、调用工具或访问外部知识，将部分任务交给更适合的执行模块完成。

从技术路线来看，大语言模型的推理增强大致经历了提示增强、工具增强以及过程控制与系统协同三个阶段。早期方法主要通过提示设计和采样策略，引导模型展开中间推理过程；随后，研究重点逐渐转向程序执行、知识检索和 API 调用，使精确计算和外部知识能够进入模型求解链路；近期工作则更关注推理过程是否可记录、可检查和可回退，尝试把规划、行动、反馈与修复组织成一个连续的求解过程\cite{wei2022chain,wang2023selfconsistency,gao2023pal,schick2023toolformer,qin2023toolllm,yao2023react,critic2024}。因此，复杂任务能力的提升并不只取决于模型生成更长的推理文本，还取决于模型能否在不同阶段完成合适的子任务，并与外部执行能力形成稳定配合。

需要注意的是，上述推理增强方法多数面向自然语言问答、数学推理或通用代码生成任务。此类任务的输入通常可以被处理为线性文本，输出形式也多为答案、程序或行动序列。图推理任务则不同，其输入包含节点、边、属性和约束条件，求解过程还需要满足明确的拓扑关系。若仅扩展一般思维链，模型可能在中间步骤中遗漏节点关系、边属性或约束条件；若只依赖工具调用，也可能出现检索到的知识、选择的算法和最终程序实现不一致的问题。因此，已有推理增强研究虽然为本文提供了分阶段求解、工具调用和过程验证方面的方法基础，但在高约束图推理场景中，还需要进一步引入面向图结构的中间表示和可执行验证机制\cite{besta2024graph,gao2023pal,yao2023react,lightman2023processreward}。

\subsection{面向图推理任务的大语言模型研究}

针对图推理问题，近年来出现了一批专门的评测基准。NLGraph 关注自然语言条件下的经典图算法任务\cite{chen2024nlgraph}，GraphWiz 面向指令式图算法推理\cite{gong2024graphwiz}，GraphArena、GraphInstruct 与 LLM4DyG 则分别从经典图问题、指令跟随场景和动态图场景评估模型能力\cite{tang2025grapharena,graphinstruct2024,llm4dyg2024}。相关方法研究也开始从图理解对齐、图序列化表达和组合优化思维生成等角度系统推动该方向发展\cite{fatemi2024talk,zhang2023graphmeetsllms,jin2024llmgraphs,gundam2024,graphthought2025,pseudocode2025}。这些基准与方法研究普遍表明，尽管大模型在小规模样例和规则明确的问题上能够取得一定效果，但在结构复杂、约束严格或需要组合优化的场景下，其性能仍显著受限。

现有图推理增强研究可以概括为三种思路。第一种思路侧重于增强模型的图结构感知能力，通常通过任务改写、结构化提示或面向图任务的指令微调，使模型更容易识别节点、边及其拓扑关系\cite{tang2024graphgpt,chai2023graphllm}。第二种思路侧重于程序化求解，即把自然语言中的图任务转换为 Python 程序，再借助 NetworkX 等图计算库完成具体计算\cite{hagberg2008networkx,zhang2025gcoder}。第三种思路则引入多智能体协作机制，让不同角色分别负责规划、检索、代码生成、评审和修复等环节，以减少单一模型一次性完成全部步骤时产生的错误\cite{li2024graphteam}。

不同图推理基准测试集在能力考察上并不完全相同。NLGraph 主要考察自然语言描述到经典图算法之间的对应关系，重点在于模型能否从文本中识别任务语义并选择合适的图算法。GraphWiz 和 GraphInstruct 更关注指令跟随能力、推理步骤的一致性以及最终输出格式是否符合要求。GraphArena 的任务范围更广，也包含更多组合优化类问题；LLM4DyG 则进一步引入时间演化和动态图约束，用于考察模型对动态关系的处理能力\cite{chen2024nlgraph,gong2024graphwiz,tang2025grapharena,graphinstruct2024,llm4dyg2024}。近期围绕图推理泛化、复杂图理解、大规模图数据编码和强化学习式图推理的研究也表明，图推理评测不能只看最终答案是否正确，还需要同时关注输入解析、图构建、算法选择、代码执行和结果验证等环节\cite{zhang2024memorization,yuan2024gracore,wu2025grapheval36k,guo2025g1}。因此，图推理更接近一个由多个子过程共同组成的系统性任务，而不是单一的文本推理能力测试。

从已有工作的演进过程看，图推理增强方法已经逐渐形成结构感知增强、程序化求解和多智能体协作三条路线。结构感知类方法主要试图缓解语言模型对图结构表示不足的问题；程序化方法通过调用成熟图计算库来提高结果的可执行性和确定性；多智能体方法则把任务理解、规划组织、知识获取和错误修复拆分给不同角色完成\cite{tang2024graphgpt,chai2023graphllm,zhang2025gcoder,li2024graphteam}。这些研究说明，图推理研究的重点正在从单纯提升模型的结构理解能力，转向构建更加稳定的图任务求解流程。不过，在中间规划表示如何构造、检索到的知识如何与算法步骤保持一致，以及执行错误如何形成有效反馈闭环等方面，现有方法仍有进一步改进空间。这些问题也是下一小节需要继续讨论的重点。

\subsection{研究现状总结}

尽管已有研究取得了积极进展，但从图推理任务的实际求解链路看，现有方法尚未充分解决一个核心问题：如何在任务理解、知识调用与程序执行之间建立稳定且可验证的中间组织机制。综合前述研究可以看到，当前方法的不足主要集中在两条相互关联的技术路线上。

第一类问题主要体现在端到端图推理链路中。此类方法通常要求模型在一次生成过程中同时完成任务语义理解、图结构恢复、算法选择、外部知识调用和结果整理。由于这些步骤没有被明确拆分，任一环节中的偏差都可能继续传递到后续结果中。当图结构被序列化为较长文本，或任务同时包含属性条件、路径约束和组合优化目标时，模型更容易遗漏节点与边之间的局部关系，进而产生约束不满足或答案幻觉等问题。即使系统接入外部文档或 API，如果检索查询仍直接由原始问题构造，而缺少对算法步骤的中间规划约束，返回的知识也未必能够与实际求解过程保持一致。

第二类问题主要来自对历史经验数据和专门适配过程的依赖。部分程序生成或经验增强方法通过任务--代码对、样例库检索或特定任务微调来改善图推理表现。这类方法在任务类型和输入形式较为固定时通常有效，但其效果往往依赖预先构建的样例、代码经验和适配数据。随着任务类型变化、图规模扩大或测试基准迁移，原有经验库可能难以覆盖新的求解场景，同时还会带来数据构建、维护和更新方面的额外成本。因此，这类方法在跨任务泛化和长期适配方面仍然存在一定限制。

因此，有必要构建一种更加稳健的图推理框架，在保留大语言模型通用理解能力的同时，将规划、知识使用与执行验证显式拆分，以提升求解过程的可靠性与泛化性能。

\section{研究内容与研究目标}

围绕上述问题，本文提出一种基于规划反馈的图推理多智能体协同求解框架 PlanGraph。本文所关注的研究对象，不是一般意义上的文本推理增强，而是如何让大语言模型在图推理任务中形成可规划、可验证、可修正的求解过程。结合前一节的分析，本文认为改进空间主要体现在两个方面：其一，通过显式规划对复杂图任务进行结构化拆解，减轻模型在一次生成中同时处理任务理解、算法匹配与约束满足的负担；其二，以规划结果而非原始问题作为知识调用的中间语义，使检索目标更贴近当前求解步骤，并在必要时触发知识片段精炼或补充检索。基于这一定位，PlanGraph 不直接从原始自然语言问题生成最终程序或答案，而是将整个求解过程拆分为规划生成、规划引导检索与检索质量评估、执行验证与反馈推理三个阶段。

在规划生成阶段，系统首先从原始问题中识别任务目标、输入图信息以及最终输出要求，并由规划智能体生成结构化伪代码。这里的伪代码并不是完整程序，而是对算法意图、求解顺序和关键约束的概括，用来描述后续代码生成需要遵循的主要操作。进入知识检索阶段后，系统不再直接使用原始问题文本作为查询，而是依据规划阶段得到的伪代码，从图算法文档知识库中检索相关 API 说明、输入输出约束和使用示例。这样可以使检索内容更接近实际求解步骤，减少与当前任务无关的知识片段对代码生成的干扰。在执行验证与反馈推理阶段，编码智能体结合规划结果和检索知识生成可执行代码；系统随后在小样本图上进行单元测试，用于检查程序的基本正确性。若执行过程中出现报错、输出格式异常或结果不满足约束，推理智能体会根据反馈信息进行修正；若多轮局部修正仍无法通过验证，系统则触发重新生成机制。通过这种局部修正与重生成结合的双层闭环，规划在前向过程中约束检索和代码生成，执行反馈则在反向过程中修正实现偏差、知识不匹配和约束遗漏，最终得到更可靠的答案。

基于上述分析，本文将研究问题归纳为两个相互关联的挑战。第一，如何将复杂的自然语言图推理任务稳定转换为可执行的求解过程。图任务通常同时包含任务语义、图结构信息和约束条件。如果缺少合适的中间表示，模型可能出现“理解了问题但程序实现偏离目标”，或“程序可以运行但遗漏关键约束”的情况。第二，如何在规划约束下完成面向当前任务的外部知识调用，并进一步提高程序运行的可靠性。对于图推理任务而言，检索到的 API、算法说明和示例是否与当前规划步骤一致，会直接影响后续代码生成质量。同时，单次程序生成很难覆盖所有边界情况，因此还需要通过执行反馈和迭代修复机制来提高最终答案的正确率。

围绕上述问题，本文的总体研究目标是构建一种兼顾可解释性、稳定性和泛化能力的图推理框架。该框架不仅要求大语言模型给出最终答案，还要求其形成可追踪、可验证并且可以修正的求解过程。具体而言，本文希望实现以下目标：

（1）在方法层面，建立规划生成、知识检索和验证推理相结合的分阶段求解范式，并通过反馈回流减少推理过程中的错误传播和答案幻觉；

（2）在协同推理机制层面，设计面向多智能体协同的共享状态维护机制和黑板调度机制，使不同模块之间能够保持信息一致，并提高求解过程的可解释性；

（3）在实验层面，基于多个公开图推理基准测试集评估所提方法，重点考察其在精确匹配率、稳定性和鲁棒性等方面的表现。

在上述研究内容与研究目标基础上，本文形成如下技术路线。整体求解过程可概括为问题解析—显式规划—知识检索—程序生成—执行验证—答案组织六个阶段。前文所述的规划生成、规划引导检索与验证推理是方法层面的抽象划分，而这里的六个阶段是系统执行层面对该方法的细化。如图\ref{fig:research_roadmap}所示，本文整体研究思路可以归纳为两条互相衔接的主线：一条主线关注任务理解与规划建模，通过问题智能体和规划智能体把自然语言图任务转化为稳定的中间伪代码表示；另一条主线关注知识增强与程序化求解，通过检索、编码和验证闭环把规划结果落实为可执行且可复核的求解过程。

\begin{figure}[htbp]
\centering
\includegraphics[width=0.90\textwidth]{figures/paper/research_roadmap.pdf}
\caption{研究技术路线图}
\label{fig:research_roadmap}
\end{figure}

系统首先对输入任务进行结构化解析，提取任务类型、图中涉及的节点与边信息以及相应约束条件。随后，规划模块将求解思路整理为算法级伪代码，作为检索和代码生成阶段共享的中间语义表示。检索模块根据伪代码中的算法步骤，从外部文档知识库中筛选可能需要的 API 说明和算法描述，并结合检索结果的匹配情况判断是否需要进一步精炼知识片段或发起补充检索。之后，编码模块依据规划结果和检索知识生成可执行程序，验证模块则通过测试结果判断程序是否满足任务要求，并在出现错误时触发修复或重新生成。最后，系统按照题目要求整理并输出可复核的答案。相比直接生成最终答案的方式，该技术路线把复杂图任务拆分为多个可以单独检查的环节，使错误来源更容易被定位；同时，执行反馈为后续修正提供了明确依据，避免系统在失败后只是重复生成缺乏针对性的答案。


围绕上述研究目标与技术路线，本文的主要贡献可概括为以下三个方面：



（1）提出了一种面向图推理任务的规划反馈协同求解方法。针对现有大语言模型图推理中任务理解、知识调用和程序执行高度耦合的问题，本文在自然语言问题与程序生成之间引入结构化规划层，并以规划结果约束知识检索和代码生成过程。相比直接提示、一般思维链推理或单纯程序生成方法，该方法将图推理组织为“规划约束生成—执行反馈修正”的闭环过程，使图结构约束、算法步骤和输出要求能够在求解链路中被显式保留和检查。

（2）构建了一个基于黑板调度的多智能体图推理协同框架。该框架围绕问题解析、规划生成、知识检索、代码生成、执行验证和反馈修复设置不同功能智能体，并通过共享黑板、共享缓存链和主流程调度机制实现跨阶段协同。其运行方式不是依赖智能体之间的自由对话来推进任务，而是在统一控制流下记录阶段结果、执行状态和恢复记录，从而增强复杂图任务求解过程的可追踪性、可恢复性和可复现性。

（3）在多个公开图推理基准测试集上验证了 PlanGraph 的有效性和适用边界。实验覆盖五个公开图推理基准测试集，从总体性能、任务类别、NP 完全任务、样本难度、检索质量、参数敏感性、鲁棒性和模块消融等角度进行分析。

\section{论文组织结构}

全文结构安排如下。

（1）第\ref{chp:introduction}章为绪论，主要介绍研究背景与意义、国内外研究现状，以及本文的研究内容、研究目标和技术路线。

（2）第\ref{chp:related_tech}章梳理与本文相关的理论基础和关键技术，包括图推理任务基础、大语言模型推理增强、工具增强与多智能体协同，以及相关技术总结等内容。

（3）第\ref{chp:method}章重点介绍基于规划反馈的图推理求解方法，围绕研究问题定义、基于结构化任务的规划生成方法、基于伪代码规划的知识检索方法、基于执行反馈的推理验证方法、方法有效性分析以及方法验证展开。

（4）第\ref{chp:system}章介绍基于黑板调度的多智能体协同推理机制，重点说明研究问题定义、共享黑板状态建模、错误归因调度、共享缓存反馈恢复、运行记录诊断以及协同推理机制验证。

（5）第\ref{chp:experiment}章给出综合实验验证，在前两章分模块验证的基础上，统一说明实验设置、总体性能结果和性能收益来源，并对 PlanGraph 的整体有效性、失败样例和适用边界进行综合分析。

（6）第\ref{chp:conclusion}章对全文工作进行总结，概括本文的主要结论，并对后续可能的改进方向和研究展望进行讨论。
